% chapter_06_Decay.tex -- Chapter 06, Decay: the hadron decays, the full generator event.
% Every number comes from test/06_Decay/out/ (ledger_06_Decay.txt, genPartAnalyzer_06_Decay.log,
% step1_GEN_06_Decay.log) and from FWLite reads of step1_GEN_06_Decay.root listed in chapter_06_Decay.md.
\chapter{06 -- Decay: the primary hadrons decay and the generator event is complete}
\label{ch:06_Decay}

\section{What this chapter is about}
Chapter~\ref{ch:05_Had} left the event as 320 primary hadrons: the mesons and baryons that the strings of
the hadronisation produced, among them the $B^-$ that carries the $b$ quark and the $B^{*+}$ that carries the
$\bar{b}$. Most of these hadrons do not live long enough to reach the detector. This chapter switches on the
last step of Pythia8, \code{HadronLevel:Decay = on}, and lets them decay: the $\pi^0$ into two photons, the
$\rho$ into pions, the $D$ and $B$ mesons through their weak decay chains, down to the particles that Pythia8
calls stable. The file written here, \code{test/06\_Decay/out/step1\_GEN\_06\_Decay.root}, is the complete
generator event: its 634 stable particles are what the jet clustering of chapter~\ref{ch:07_GenJet} reads, what
the pileup of chapter~\ref{ch:08_PU} is added to, and what chapter~\ref{ch:09_G4} hands to Geant4. The book
follows this file from here on; every later chapter starts from it.

The chapter answers three questions. Which hadrons decay in Pythia8 and which are left, still undecayed, to
Geant4? What does the decay of the two $b$ hadrons of the event look like, and where does the energy go,
in particular the neutrino that the detector will never see? And what is the final state made of, by kind and
by region, once everything that can decay has decayed? The ``Open problems'' section then looks at what the
release uses for the decay tables and the lifetime boundary, and at the age of the common Pythia8 settings.

\begin{plainwords}[Stable is a decision, not a fact]
A particle that lives $10^{-10}$~s before decaying travels centimetres in the detector; a particle that lives
$10^{-16}$~s decays before it has moved a micron. The generator and the detector simulation share the work: Pythia8
decays the short-lived ones, Geant4 decays the long-lived ones on their way through the material, because
for those it matters where they decay. The boundary is a number in the configuration, a proper lifetime
$c\tau_0$ of $10\mm$. ``Status~1'' (stable) in the generator record means ``left for Geant4'', not
``lives forever''.
\end{plainwords}

\section{From the previous chapter}
The parton level and the hadronisation are identical to chapter~\ref{ch:05_Had}: same fragment, same seeds,
and Pythia8 does the same 32 MPI, 28 ISR and 64 FSR emissions and the same string fragmentation
before the new switch acts.\footnote{\path{test/06_Decay/out/step1_GEN_06_Decay.log}, the Info listing; the
record by status group in \path{ledger_06_Decay.txt} (4, 124, 190, 192, 233, 165 entries for the hard process,
MPI, ISR, FSR, beam remnants and strings) is the same as in \path{test/05_Had/out/ledger_05_Had.txt}.} The
analyzer log confirms it stage by stage: the $B^-$ at $\pt = 94.28\GeV$, $\eta = -1.57$ and the $B^{*+}$ at
$\pt = 114.17\GeV$, $\eta = -3.75$ are the same primary hadrons, at the same place, as in
chapter~\ref{ch:05_Had}, and the 320 particles of stage~4 of this chapter are the 320 status-1 particles of the
previous one.\footnote{\path{test/06_Decay/out/genPartAnalyzer_06_Decay.log}, the lines \code{stage 4} and
\code{stage 5}.}

What changes is the end of the record. In chapter~\ref{ch:05_Had} all 320 primary hadrons had status~1. Here
188 of them decay and receive status~2 (``decayed''), their decay products are appended to the record, and
those products that are themselves unstable decay in turn (114 further decays), until only status-1 particles
are left. The record grows from 1230 to 1846 entries: 616 new particles, of which 502 are stable and 114 decayed
again, so the event ends with 634 stable particles and 302 decayed ones, and 132 of the primary hadrons survive
untouched because their lifetime is above the boundary (table~\ref{tab:06_ledger}, and the legend of
figure~\ref{fig:06_display}: 9 + 16 + 132 + 477 = 634). The two $b$ hadrons are followed through their decays:
the 9 stable products of the $B^-$ and the 16 of the $B^{*+}$ still carry exactly the \pt{} of their parents,
94.28 and $114.17\GeV$, and the net $\pt$ of $b + \bar{b}$ stays $20.30\GeV$.

\section{The numbers of the event}
\paragraph{Who decays in Pythia8 and who is left to Geant4.}
Two lines of the CMS common settings decide it:\footnote{\path{Configuration/Generator/python/Pythia8CommonSettings_cfi.py}
of the release, copied into \path{test/06_Decay/step1_GEN_06_Decay_cfg.py} as the block \code{pythia8CommonSettings}.}
\code{ParticleDecays:limitTau0 = on} and \code{ParticleDecays:tau0Max = 10}: a particle species is decayed by
Pythia8 only if its proper lifetime $c\tau_0$, as listed in Pythia8's particle data, is below $10\mm$. The
particle table of Pythia 8.311\footnote{\path{share/Pythia8/xmldoc/ParticleData.xml} of the Pythia8 of the release,
\code{tau0} in mm; the GEN log prints ``no particle data has been changed from its default value''.} then sorts the
hadrons of the event into the two groups of table~\ref{tab:06_who}. The $\pi^0$ ($c\tau_0 = 2.6 \times 10^{-5}\mm$),
the $\eta$, the strong resonances ($\rho$, $\omega$, $K^*$, $\Delta$, which have a width and no lifetime at all) and
the weakly decaying charm and bottom hadrons ($D^0$: $0.123\mm$, $B^+$: $0.491\mm$) decay here. The $K^0_S$
($26.8\mm$), the $\Lambda$ ($78.9\mm$), the $\Sigma^\pm$, $\Xi$, $\Omega$ (24 to 87~mm), the charged pions
and kaons (7.8 and 3.7~m), the $K^0_L$ (15.3~m), the neutron and the muon are left with status~1: Geant4 will
decay them, if they decay at all before being absorbed.

\begin{table}[htbp]
\centering\small
\caption{The hadron decays of the event: what Pythia8 decayed (status 2, 302 decays) and what it left stable
(status 1, 634 particles), by species, with the proper lifetime $c\tau_0$ of Pythia8's particle data that decides
the group. From an FWLite read of \code{test/06\_Decay/out/step1\_GEN\_06\_Decay.root} (listed in the companion) and
\code{ParticleData.xml} of Pythia~8.311.}
\label{tab:06_who}
\begin{tabular}{@{}llrl@{}}
\toprule
 & species & $n$ & $c\tau_0$ (Pythia8) \\
\midrule
\multirow{8}{*}{\shortstack[l]{decayed by Pythia8\\(status 2)}}
 & $\pi^0$ & 131 & $2.6\times10^{-5}\mm$ \\
 & $\rho^0$, $\rho^\pm$ & 60 & resonance, $\Gamma = 149\MeV$ \\
 & $K^0$, $\bar{K}^0$ (to $K^0_S$ / $K^0_L$) & 33 & $10^{-6}\mm$ (flavour eigenstate) \\
 & $K^{*}$, $\omega$, $\eta$, $\eta'$, $\phi$ & 58 & resonances; $\eta$: $1.5\times10^{-7}\mm$ \\
 & $\Delta$, $\Sigma^*$, $\Sigma^0$ & 10 & resonances; $\Sigma^0$: $2.2\times10^{-8}\mm$ \\
 & $D^0$, $\bar{D}^0$, $D^{*}$, $D^*_0$ & 7 & $D^0$: $0.123\mm$ \\
 & $B^-$, $B^+$, $B^{*+}$ & 3 & $B^\pm$: $0.491\mm$ \\
\midrule
\multirow{9}{*}{\shortstack[l]{left stable for Geant4\\(status 1)}}
 & $\gamma$ & 273 & stable \\
 & $\pi^+$, $\pi^-$ & 261 & $7805\mm$ \\
 & $K^+$, $K^-$ & 29 & $3713\mm$ \\
 & $K^0_L$ & 22 & $15\,330\mm$ \\
 & $K^0_S$ & 13 & $26.8\mm$ \\
 & $n$, $\bar{n}$; $p$, $\bar{p}$ & 13; 11 & stable in practice \\
 & $\Lambda$, $\bar{\Lambda}$; $\bar{\Sigma}^+$ & 8; 2 & $78.9\mm$; $44.3\mm$ \\
 & $e^-$; $\bar{\nu}_e$ & 1; 1 & stable; invisible \\
\bottomrule
\end{tabular}
\end{table}

\paragraph{The status codes.}
In the \code{genParticles} collection every entry carries the Pythia8 status: the codes 21 to 79 of the earlier
chapters tell in which step a parton was produced, and the two codes of this chapter are 2 (``decayed'': the
particle existed, its daughters follow) and 1 (``stable'': no daughters, the end of the generator's work). The
decay tree is stored as mother and daughter links, so the full history of every stable particle, from the
hard process through strings and decays, can be walked back in the file. The analyzer of the book draws stage~5 as
the set of status-1 particles, coloured by origin, and the $b$ and $\bar{b}$ systems as the stable descendants of
the two primary $b$ hadrons.

\paragraph{The decay of the $b$ hadrons.}
The $B^-$ (index 829 in \code{genParticles}) decays semileptonically: $B^- \to e^-\,\bar{\nu}_e\,D^{*0}_0(2300)$,
followed by $D^{*0}_0 \to D^0\pi^0$, $D^0 \to K^{*-}\rho^+$, $K^{*-} \to \bar{K}^0\pi^-$ with $\bar{K}^0 \to K^0_S$,
$\rho^+ \to \pi^+\pi^0$ and two $\pi^0 \to \gamma\gamma$. Its 9 stable products are the antineutrino
($E = 95.5\GeV$, $\pt = 37.5\GeV$, $\eta = -1.59$), the electron ($10.9\GeV$), a $K^0_S$ ($20.6\GeV$), a $\pi^-$
($26.6\GeV$), a $\pi^+$ ($32.6\GeV$) and four photons ($44.3$, $4.1$, $0.8$ and $0.4\GeV$), and they sum to the
$235.9\GeV$ of the $B^-$. The neutrino alone takes 40\% of the $B^-$ energy and 40\% of its \pt. It is the
only neutrino of the event, and it is exactly the generator missing transverse momentum stored in the file:
\code{genMetTrue} has $\pt = 37.46\GeV$ at $\phi = -2.10$, the neutrino's direction. This is the energy that
chapter~\ref{ch:09_G4} loses at the hand-over: Geant4 is not given neutrinos, and the $b$ jet of
chapter~\ref{ch:07_GenJet} will be measured without it.

The $B^{*+}$ (index 968) first radiates a $33.2\GeV$ photon, $B^{*+} \to B^+\gamma$ (an electromagnetic decay, at the
production vertex), and the $B^+$ decays hadronically into five bodies, $B^+ \to \bar{D}^{*0}\,\pi^+\,\pi^0\,K^0\,\bar{K}^{*0}$,
with $\bar{D}^{*0} \to \bar{D}^0\gamma$ and $\bar{D}^0 \to \bar{K}^0\pi^-\pi^+\pi^0\pi^0$. Its 16 stable products
are 10 photons, two $K^0_S$, one $K^0_L$, two $\pi^+$ and one $\pi^-$, with the full $2421.4\GeV$ of the
$B^{*+}$ and no neutrino: everything of the $\bar{b}$ system is in principle visible, but at $\eta = -3.75$ it is
all in the forward calorimeter.

\paragraph{Where the decays happen.}
The $10\mm$ boundary is on the lifetime of the species, not on the distance a given particle travels. With
$\gamma\beta = p/m = 44.7$ the $B^-$ flies $7.8\mm$ before it decays; the forward $B^+$, with $p = 2388\GeV$
and $\gamma\beta = 452$, flies $117.3\mm$, and its $\bar{D}^0$ another $140.7\mm$, so that the $\bar{D}^0$ decay
vertex sits $258\mm$ from the primary vertex, at $r = 12.1\mm$ and $z = -227\mm$.\footnote{FWLite, the production
vertices of the daughters; the primary vertex of the event is at $x = 0.28$, $y = -0.17$, $z = 31.1\mm$ after the
beam-spot smearing of \code{VtxSmearedRealistic}. Of the 302 decays, 17 have a flight path above $0.01\mm$,
4 above $1\mm$ and 2 above $10\mm$.} These displaced vertices are the signature by which $b$ jets are tagged;
they are in the file, and chapter~\ref{ch:09_G4} shows which of them Geant4 accepts as production points.

\paragraph{The composition of the final state.}
Photons are the largest group by number, 273 of the 634 stable particles, and 262 of them come from the 131
$\pi^0$ (8 more from $\eta$ decays, one each from the $B^{*+}$, a $\bar{D}^{*0}$ and a $\bar{\Sigma}^0$); they
carry $2876\GeV$, of which $2703\GeV$ from the $\pi^0$. Charged hadrons are 303 (261 charged pions, 29
charged kaons, 11 protons and antiprotons, two $\bar{\Sigma}^+$) with $7168\GeV$; neutral hadrons are only 56
(22 $K^0_L$, 13 $K^0_S$, 13 neutrons, 8 $\Lambda$) with $3450\GeV$; and there is one electron and one
neutrino.\footnote{\path{test/06_Decay/out/ledger_06_Decay.txt}, ``by kind'', and the FWLite species counts of the
companion.} In chapter~\ref{ch:05_Had} the 320 primary hadrons were 169 charged and 151 neutral: the decays turn
most of the neutral hadrons (all the $\pi^0$, $\eta$, $\omega$, $K^*$) into photons and charged pions, which is why
a jet is, by number, mostly photons and charged pions, and why the electromagnetic calorimeter sees a sizeable
part of the energy even of a hadronic jet (here the photons carry 21\% of the energy of the stable particles).
Figure~\ref{fig:06_species} shows the stable species.

\begin{figure}[htbp]
\centering
\begin{tikzpicture}[x=1.05cm, y=0.0115cm]
\draw[->] (0,0) -- (11.4,0);
\draw[->] (0,0) -- (0,300);
\foreach \y in {0,100,200,300} {\draw (0.1,\y) -- (-0.1,\y) node[left,font=\tiny] {\y};}
\histbar{0}{273} \histbar{1}{261} \histbar{2}{29} \histbar{3}{22} \histbar{4}{13} \histbar{5}{13}
\histbar{6}{11} \histbar{7}{8} \histbar{8}{2} \histbar{9}{1} \histbar{10}{1}
\foreach \x/\l in {0/$\gamma$, 1/$\pi^\pm$, 2/$K^\pm$, 3/$K^0_L$, 4/$K^0_S$, 5/$n\,\bar{n}$, 6/$p\,\bar{p}$, 7/$\Lambda\,\bar\Lambda$, 8/$\bar\Sigma^+$, 9/$e^-$, 10/$\bar\nu_e$}
  {\node[below, font=\scriptsize] at (\x+0.5,0) {\l};}
\node[font=\scriptsize, anchor=north east] at (11.3,290) {634 stable particles of the event, by species};
\end{tikzpicture}
\caption{The stable (status 1) particles of the event by species, from the FWLite count of
\code{test/06\_Decay/out/step1\_GEN\_06\_Decay.root} (companion, section 4). Photons, almost all from $\pi^0$
decays, and charged pions make 84\% of the final state by number.}
\label{fig:06_species}
\end{figure}

\begin{plainwords}[Why a jet is mostly photons and charged pions]
Strings fragment into pions in roughly equal numbers of $\pi^+$, $\pi^-$ and $\pi^0$ (here 261 charged pions
against 131 $\pi^0$, with a few from decays of heavier hadrons), plus a minority of kaons and baryons. The $\pi^0$
decays at once into two photons. So by the time Geant4 starts, a jet is a bundle of charged pions, which leave
tracks and then shower in the hadron calorimeter, and photons, which shower in the electromagnetic calorimeter.
That is the basis of particle-flow reconstruction and the reason the two calorimeters share the jet energy.
\end{plainwords}

\section{Ledger and display}
Table~\ref{tab:06_ledger} closes the ledger of the stage. Decays conserve energy and momentum, so the total energy
of the stable particles stays $13\,600.0\GeV$, the two proton energies. The scalar sum of \pt{} rises from
$505.7$ to $528.9\GeV$, because a decay shares the momentum of one particle among several that are not
collinear, and for the same reason energy migrates between the $\eta$ regions: the tracker region
$\abseta < 2.5$ gains $49\GeV$ and the region beyond $\abseta = 5.5$ loses $193\GeV$ as forward hadrons decay
into products spread over a wider range of $\eta$. The visible energy within $\abseta < 5.2$ grows from $5779$
to $5851\GeV$, while $95.5\GeV$ become invisible in the one neutrino: at this stage the only hole in the event
is the semileptonic $B^-$ decay.

\begin{table}[htbp]
\centering\small
\caption{Ledger of the decay stage: the stable (status 1) particles before (chapter~\ref{ch:05_Had},
\code{test/05\_Had/out/ledger\_05\_Had.txt}) and after the hadron decays (\code{test/06\_Decay/out/ledger\_06\_Decay.txt}).
$n$ is the number of particles, $E$ the energy and \pt{} the scalar sum of transverse momenta, in GeV.
The rows by region count every stable particle, the neutrino (at $\eta=-1.59$) included.}
\label{tab:06_ledger}
\begin{tabular}{@{}lrrrrrrrrr@{}}
\toprule
 & \multicolumn{3}{c}{before (05\_Had)} & \multicolumn{3}{c}{after (06\_Decay)} & \multicolumn{3}{c}{difference} \\
\cmidrule(lr){2-4}\cmidrule(lr){5-7}\cmidrule(lr){8-10}
 & $n$ & $E$ & \pt & $n$ & $E$ & \pt & $n$ & $E$ & \pt \\
\midrule
entries in \code{genParticles} & 1230 & & & 1846 & & & $+616$ & & \\
decayed (status 2)             & 0 & & & 302 & & & $+302$ & & \\
stable (status 1)              & 320 & 13600.0 & 505.7 & 634 & 13600.0 & 528.9 & $+314$ & 0.0 & $+23.2$ \\
\midrule
charged hadrons & 169 & 8863.8 & 355.0 & 303 & 7167.7 & 259.1 & $+134$ & $-1696.1$ & $-95.9$ \\
neutral hadrons & 151 & 4736.2 & 150.7 &  56 & 3449.9 &  84.7 & $-95$  & $-1286.3$ & $-66.0$ \\
photons         &   0 &    0.0 &   0.0 & 273 & 2875.9 & 142.7 & $+273$ & $+2875.9$ & $+142.7$ \\
electrons       &   0 &    0.0 &   0.0 &   1 &   10.9 &   5.0 & $+1$   & $+10.9$   & $+5.0$ \\
neutrinos       &   0 &    0.0 &   0.0 &   1 &   95.5 &  37.5 & $+1$   & $+95.5$   & $+37.5$ \\
\midrule
$\abseta < 2.5$ (tracker)        & 141 &  684.6 & 259.9 & 316 &  733.8 & 278.7 & $+175$ & $+49.2$  & $+18.8$ \\
$2.5 \le \abseta < 3.0$          &  30 &  196.8 &  24.6 &  55 &  193.1 &  24.6 & $+25$  & $-3.7$   & $0.0$ \\
$3.0 \le \abseta < 5.2$ (HF)     &  94 & 4897.9 & 192.4 & 189 & 5019.6 & 197.4 & $+95$  & $+121.7$ & $+5.0$ \\
$5.2 \le \abseta < 5.5$          &  15 &  810.7 &   7.9 &  16 &  836.4 &   7.8 & $+1$   & $+25.7$  & $-0.1$ \\
$\abseta \ge 5.5$ (not simulated) & 40 & 7010.1 &  20.9 &  58 & 6817.1 &  20.3 & $+18$  & $-193.0$ & $-0.6$ \\
\midrule
visible ($\abseta < 5.2$, no $\nu$) & 265 & 5779.2 & 476.9 & 559 & 5851.0 & 463.4 & $+294$ & $+71.8$ & $-13.5$ \\
invisible (neutrinos)               &   0 &    0.0 &   0.0 &   1 &   95.5 &  37.5 & $+1$   & $+95.5$ & $+37.5$ \\
\bottomrule
\end{tabular}
\end{table}

Figure~\ref{fig:06_display} is the $\eta$--$\phi$ display of the complete event. Compared with the stage-4
display of chapter~\ref{ch:05_Had}, where the $B^-$ and the $B^{*+}$ were two single large markers, the two
$b$ systems are now clusters of 9 (red) and 16 (blue) small markers around the same positions, $(\eta, \phi) =
(-1.57, -2.09)$ and $(-3.75, 1.01)$: the decay products of a $236\GeV$ and a $2421\GeV$ hadron stay within a
few tenths of a unit of their parent. This is the picture a jet algorithm will see in the next chapter. The
grey points are the 477 other decay products, the violet ones the 132 primary hadrons that did not decay; 38
particles beyond $\abseta = 6$ are drawn on the border.

\begin{figure}[htbp]
\centering
\includegraphics[width=\textwidth]{06_Decay/out/genPartAnalyzer_06_Decay_page-5.png}
\caption{The $\eta$--$\phi$ display of the complete generator event, stage 5 of
\code{test/06\_Decay/out/genPartAnalyzer\_06\_Decay.pdf} (page 5): the 634 stable particles, marker size growing with
\pt, coloured by origin. Red: the 9 stable decay products of the $B^-$; blue: the 16 of the $B^{*+}$; violet: the
primary hadrons left undecayed; grey: the decay products of the other hadrons. The magenta star is the vector sum
of the $b$ and $\bar{b}$ systems, $\pt = 20.3\GeV$.}
\label{fig:06_display}
\end{figure}

\FloatBarrier
\section{Open problems for Phase-2}
\paragraph{Which decay tables.}
This event is decayed with Pythia8's own tables: the GEN log prints ``no particle data has been changed from its
default value'', and the configuration has no \code{ExternalDecays} block. That is also what the Phase-2 release
validation does: the fragments it uses, \path{Configuration/Generator/python/TTbar_14TeV_TuneCP5_cfi.py} and
\path{MinBias_14TeV_pythia8_TuneCP5_cfi.py}, contain no external decayer either. The production fragments that
do use EvtGen do it through the \code{ExternalDecays} parameter set of the Pythia8 module, for example in
\path{Configuration/Generator/python/Hadronizer_TuneCP5_13TeV_aMCatNLO_FXFX_5f_max2j_max1p_LHE_pythia8_evtgen_cff.py}:
\code{EvtGen130} with \code{decay\_table} set to \path{GeneratorInterface/EvtGenInterface/data/DECAY_2010.DEC},
\code{particle\_property\_file} to \code{evt.pdl}, and empty \code{list\_forced\_decays} and \code{operates\_on\_particles}
(so EvtGen takes every species it knows, $b$ and $c$ hadrons among them). The $B$-physics fragments of the same
directory use \code{DECAY\_2014\_NOLONGLIFE.DEC} with \code{evt\_2014.pdl}, while the data package shipped with
the release (\path{cms-data/GeneratorInterface-EvtGenInterface}, V02-12-00, EvtGen~2.0.0) also carries
\code{DECAY\_2020.DEC} and \code{evt\_2020.pdl}, whose $B^-$ block still reads ``Updated to PDG 2008'' with
branching fractions marked PDG2014. Three generations of tables coexist in one release, and a 2010 table is
the default of the standard hadronizer fragment. The central Phase-2 Monte Carlo fragments are not part of the
release (they live in the production repository), so which table they use could not be established here;
what the release shows is that nothing selects a newer table by default, and that the Pythia8-only path of this
chapter is the one used for the Phase-2 relvals. For a $b$-jet study at the HL-LHC the semileptonic branching
fractions and the $D$ decay multiplicities, which set the neutrino energy and the track multiplicity in the jet,
are decided by this choice.

\paragraph{The lifetime boundary and the displaced vertices.}
\code{ParticleDecays:tau0Max = 10} (mm) in \path{Pythia8CommonSettings_cfi.py} is Pythia8's own default
(\code{ParticleDecays.xml}: \code{default="10."}); CMS only switches \code{limitTau0} on. The boundary has been
the same for Run~2, Run~3 and the Phase-2 relvals. It is a cut on the species, not on the vertex: in this event
the forward $B^+$ decays $117\mm$ from the primary vertex and its $\bar{D}^0$ at $258\mm$, inside the
acceptance window that Geant4 applies to generator vertices, \code{RDecLenCut = 2.9} and \code{LDecLenCut = 30.0} (cm)
in \path{SimG4Core/Application/python/g4SimHits_cfi.py}, but a $B$ hadron of a few TeV at larger $\eta$ would
not be. Everything above the boundary ($K^0_S$, $\Lambda$, $\Sigma$, $\Xi$, $\Omega$) decays in Geant4 with
Geant4's own branching fractions and kinematics, not Pythia8's; the generator-level $V^0$ content of a jet
is therefore split between two programs. For Phase-2 the vertex smearing changes module:
\path{IOMC/EventVertexGenerators/python/VtxSmearedRealisticHLLHC_cfi.py} (\code{HLLHCEvtVtxGenerator}, read from
the \code{SimBeamSpotHLLHCObjectsRcd} of the global tag), with the crab-cavity parameters of
\code{HLLHCVtxSmearingParameters} in \path{VtxSmearedParameters_cfi.py} (\code{CrossingAngleInurad = 510},
\code{BetaCrossingPlaneInm = 0.20}, \code{EprotonInGeV = 6500} in the base set, 7000 in the 14~TeV clone), which is
what the relval scenario \code{GenSimHLBeamSpot} of \path{Configuration/PyReleaseValidation/python/upgradeWorkflowComponents.py}
selects (\code{BeamSpot: DBrealisticHLLHC}, geometry \code{ExtendedRun4D110}, era \code{Phase2C17I13M9}). The
Phase-2 inner tracker that will resolve these vertices is described in the XML of
\path{Geometry/TrackerCommonData/data/PhaseII/}; its first-layer radius is not a python parameter and was not
extracted here.

\paragraph{The age of the common settings.}
\path{Pythia8CommonSettings_cfi.py} has eight lines and none of them is tune-specific:
\code{Tune:preferLHAPDF = 2} is the Pythia8 default (\code{Tunes.xml}) and does nothing;
\code{Main:timesAllowErrors = 10000} raises the number of flawed events Pythia8 may skip before it aborts from the
default 10, so that a sick configuration is not noticed; \code{Check:epTolErr = 0.01} loosens the
energy-momentum conservation check from the default $10^{-4}$ to 1\% of the collision energy, i.e. $136\GeV$ at
13.6~TeV; \code{SLHA:minMassSM = 1000} and \code{Beams:setProductionScalesFromLHEF = off} concern SUSY and
Les Houches input; \code{ParticleDecays:allowPhotonRadiation = on} adds QED radiation in decays to lepton pairs,
which the Pythia8 manual leaves off by default to avoid double counting with an external QED program such as
PHOTOS, which EvtGen runs. The block dates from the 2017 tune campaign and its physics content is the same in
CMSSW\_15\_0\_5. The tune itself is older than the beam energy: the fragment of the book imports
\path{Configuration/Generator/python/MCTunes2017/PythiaCP5Settings_cfi.py}, whose total and elastic cross
sections are the 13~TeV values, \code{SigmaTotal:sigmaTot = 100.309} and \code{sigmaEl = 21.89} (mb), while the
release also holds \path{MCTunesRun3ECM13p6TeV/PythiaCP5Settings_cfi.py} with 101.037 and 22.08~mb for
13.6~TeV; the 14~TeV fragments of the Phase-2 relvals (\code{TTbar\_14TeV\_TuneCP5\_cfi.py},
\code{MinBias\_14TeV\_pythia8\_TuneCP5\_cfi.py}) import the 2017 block as well, so no CP5 setting has been
re-derived for 14~TeV. These cross sections feed the MPI machinery (the log of this chapter computes
$\sigma_{\mathrm{ND}} = 55.42\mb$ from them) and, through the minimum-bias events of chapter~\ref{ch:08_PU},
the number of particles per pileup collision at PU~200.

\begin{plainwords}[What to take away]
The decays are the last step of the generator and the first place where the detector enters the physics choices:
a lifetime boundary of $10\mm$ decides who decays in Pythia8 and who in Geant4, the decay tables decide how often
a $b$ hadron produces a neutrino, and this event has one, $95.5\GeV$, 40\% of the $B^-$, which no later
stage will recover. Everything downstream, jets, pileup, Geant4, digitisation, reads the 634 stable particles of
this file.
\end{plainwords}
